\section{Degeneration, exact sign transfer, and rank refinement}\label{sec:F}

We work with a complex klt germ $x\in S$ of dimension $n$, a nonzero
regular function $h$ such that $(S,H)$ is lc, where $H=\divisor(h)$,
and a rational top form $\omega$ with $\divisor_S(\omega)=0$.
Thus the reflexive canonical sheaf is trivialized by $\omega$.
A finite cyclic group $G$ fixes $x$ and acts quasi-\'{e}tale on $S$;
we replace it by its effective image.  Suppose that $h$ and $\omega$
have the same character $\chi$.  In particular, $\omega/h$ descends
to the quotient.  Fix invariant ratios
\[
 f_j=F_j/h^{m_j},\qquad F_j\in\mathcal O_{S,x}\setminus\{0\},
 \quad m_j\in\N_{>0}.
\]
The common slice consists of the $h$-lc valuations satisfying
$q(f_j)\ge0$, with the convention on centres from Section~\ref{sec:E}:
the closed image of the centre contains $x$.  Unless another
normalization is displayed, $q(h)=A_S(q)=1$.

Let $w$ be a centred $G$-invariant quasi-monomial valuation for which
\[
 \breve R=\operatorname{gr}_w\mathcal O_{S,x}
\]
is finitely generated, normal, and klt.  These are hypotheses here;
they are supplied in our applications by the exact optimizer of
Section~\ref{sec:E}.  Write $S^0=\Spec\breve R$, with vertex $o$,
and denote initial symbols by a bar.  The fine grading is by the
finitely generated value group $\Lambda_w\subset\R$, so evaluation
at $w$ is injective on this lattice.  Its torus will be denoted by
$\mathbb T_w$.  Taking invariants in each filtered exact sequence gives
\begin{equation}\label{eq:F-invariant-graded}
 \operatorname{gr}_w(\mathcal O_{S,x}^{G})=\breve R^{G}.
\end{equation}
Indeed, the Reynolds operator makes taking invariants exact in
characteristic zero.  The weight lattices of $\breve R$ and
$\breve R^G$ have finite index in one another, and hence the same rank.

\subsection{Presentations and rational approximations}

\begin{lemma}[Weighted presentations]\label{lem:F-presentation}
Choose $x_1,\ldots,x_N\in\mathfrak m_x$ whose symbols generate
$\breve R$, and put $\alpha_i=w(x_i)>0$.  Any prescribed finite
list of functions of positive $w$-value can be included in this list.
Then
\begin{equation}\label{eq:F-presentation}
 \widehat{\mathcal O}_{S,x}=\C[[X_1,\ldots,X_N]]/I,
 \qquad
 \breve R=\C[X_1,\ldots,X_N]/\operatorname{in}_{\alpha}I,
\end{equation}
where initial forms have least weight.  The valuation ideals are
exactly the images of the ideals generated by monomials of weight
at least the prescribed threshold.  The same description holds
before completion by intersection.

For every sufficiently nearby linear functional
$L:\Lambda_w\otimes\R\longrightarrow\R$, positive on the listed
weights, the vector $\alpha_L=(L(\alpha_i))$ gives the same initial
ideal, with its grading composed with $L$.  Its filtration defines
a centred quasi-monomial valuation $w_L$.  The functionals can be
chosen rational, and the resulting valuations can be computed
monomially on a fixed smooth chart for $w$.
\end{lemma}

\begin{proof}
Subduction subtracts a polynomial in the $x_i$ with the prescribed
initial symbol, then repeats at a strictly larger value.  There are
only finitely many values in a bounded interval: each is a
nonnegative integral combination of the finitely many positive
$\alpha_i$.  The Izumi upper bound for the fixed valuation $w$
implies that remainders whose $w$-values tend to infinity tend to
zero in the $\mathfrak m_x$-adic topology.  In particular
$\mathfrak m_x\subset(x_1,\ldots,x_N)+\mathfrak m_x^2$, and
Nakayama's lemma shows that the $x_i$ generate $\mathfrak m_x$.
Completed subduction proves both assertions in
\eqref{eq:F-presentation} and the assertion about threshold ideals.
Monomial threshold ideals are finitely generated.  Their completed
extensions intersect the local ring in the original ideals, by
faithful flatness of completion.

We record why the dimension and primeness arguments used below are
valid also for other positive weight vectors.  For any such vector
$\gamma$, its monomial filtration is bounded above and below by
fixed linear reindexings of the $\mathfrak m_x$-adic filtration.
Consequently its cumulative quotient lengths grow with exponent
$n$.  The associated graded algebra is
$\C[X]/\operatorname{in}_{\gamma}I$; comparison of its positive
weighted filtration with ordinary polynomial degree shows that its
growth exponent is its Krull dimension.  Thus this dimension is
$n$.  If $\operatorname{in}_{\gamma}I$ is prime, the filtration is
multiplicative and separated, and defines a valuation.  If its
value group has rational rank $r$, the degree-zero homogeneous
fraction field is its residue field and has transcendence degree
$n-r$.  This proves Abhyankar equality.  Local monomialization of
Abhyankar valuations in characteristic zero
\cite[Proposition~3.7]{JonssonMustata} then makes the
valuation quasi-monomial.

Take finitely many series in $I$ whose $\alpha$-initials generate
$\operatorname{in}_{\alpha}I$.  Their leading faces persist after
composition with $L$: ties are preserved because their fine
weights are equal, and positivity leaves only finitely many other
terms that could compete in a fixed neighbourhood of $\alpha$.
Therefore
\[
 \operatorname{in}_{\alpha}I\subset
 \operatorname{in}_{\alpha_L}I.
\]
The first ideal is prime, and both quotients have dimension $n$.
A nonzero ideal in a finitely generated domain lowers its
dimension, so this inclusion is equality.

For the final assertion, use a smooth monomial chart for $w$ at its
stratum and expand the finitely many $x_i$ over the residue
coefficient field.  Their leading faces are finite polynomials in
the stratum parameters.  The weights of these parameters belong
to $\Lambda_w$.  Composing them with nearby rational $L$ preserves
the same leading polynomials.  The resulting monomial valuation
dominates the presentation filtration.  The induced map from the
presentation graded algebra to the graded ring of this monomial
valuation is injective: the polynomial relations among these
leading polynomials are exactly the old relations, unchanged by
regrading.  Hence the two valuations agree on every function.
\end{proof}

\subsection{The canonical calculation before lc equality}

Xu--Zhuang study specialization of non-degenerate reflexive
differential forms under stable degeneration of klt singularities
\cite[Theorem~1.1 and Section~3]{XuZhuangForms}.
The calculation below tracks a divisor-free canonical generator,
including its finite-group character and any additional grading,
under the finitely generated degeneration used here.  We give this
calculation before imposing lc equality because it will also be used
to prove that equality after two gradings are compressed.

\begin{proposition}[Canonical forms under degeneration]
\label{prop:F-canonical-degeneration}
Under the hypotheses preceding Lemma~\ref{lem:F-presentation},
$S^0$ has a divisor-free reflexive canonical generator $\omega^0$.
It has the $G$-character of $\omega$.  It is homogeneous for
$\mathbb T_w$, and its weight $\kappa$ satisfies
\begin{equation}\label{eq:F-canonical-weight}
 \langle w,\kappa\rangle=A_S(w).
\end{equation}
If an additional torus acts on the original germ, preserves $w$,
and makes $\omega$ homogeneous, its character is preserved
separately.  None of these assertions assumes $A_S(w)=w(h)$.
Moreover the action of $G$ on $S^0$ is quasi-\'{e}tale.
\end{proposition}

\begin{proof}
First track discrepancy under the rational approximations of
Lemma~\ref{lem:F-presentation}.  On a fixed monomial chart with
stratum parameters $z_1,\ldots,z_r$ and separating residue
coordinates $y_1,\ldots,y_{n-r}$, write
\[
 \omega=\psi\,
 \frac{dz_1}{z_1}\wedge\cdots\wedge\frac{dz_r}{z_r}
 \wedge dy_1\wedge\cdots\wedge dy_{n-r}.
\]
The quasi-monomial discrepancy formula gives
$A_S(w)=w(\psi)\in\Lambda_w$.  Include numerators and denominators
of the rational function $\psi$ among the finite functions whose
leading faces are tracked.  The same formula on the same chart
then gives the exact identity
\begin{equation}\label{eq:F-discrepancy-approximation}
 A_S(w_L)=L(A_S(w)).
\end{equation}
In particular this is an identity of linear weights, rather than
only a limiting assertion about discrepancies.

Choose rational $L$ and multiply it by a positive integer so that
the presentation weights are positive integers.  In this paragraph
write $v$ for this integral scaling of $w_L$, and $a_i=v(x_i)$.
After shrinking to a $G$-stable affine neighbourhood, the $x_i$
have common zero set $\{x\}$.  Define the extended Rees algebra
\begin{equation}\label{eq:F-rees}
 \mathcal R=\mathcal O(S)[t,x_1t^{-a_1},\ldots,x_Nt^{-a_N}]
       \subset\mathcal O(S)[t,t^{-1}].
\end{equation}
The threshold-ideal description proves that its special fibre is
the graded algebra in \eqref{eq:F-presentation}; positive-level
quotients are supported at $x$, so the affine and local
constructions have the same special fibre.  This algebra is of
finite type and torsion free over $\C[t]$, hence flat.  Let
$\mathcal S=\Spec\mathcal R$ and $E=(t=0)$.  The divisor $E$ is
normal, integral, reduced, and Cartier, and
$\mathcal S\setminus E=S\times\mathbb G_m$.

The total space is normal.  Away from $E$ this follows from the
normality of $S$.  At points on $E$, the nonzerodivisor $t$ and
the $S_2$ property of $\mathcal R/(t)$ imply $S_2$ for
$\mathcal R$.  At the generic point of $E$, the quotient by $t$
is a field, so the total local ring is a regular discrete
valuation ring.  These observations give $R_1$ as well.

The order along $E$ restricts to $v$ on $\C(S)$ and takes value
one on $t$.  Indeed, the residues of $f t^{-v(f)}$ are the
nonzero graded symbols of $f$, and the direct grading prevents
cancellation between distinct such leading terms.  It is the
Gauss extension of $v$.  The monomial log-form computation gives
\begin{equation}\label{eq:F-rees-log-order}
 \divisor_{\mathcal S}\left(\omega\wedge\frac{dt}{t}\right)
       =(A_S(v)-1)E.
\end{equation}
For example, on a model where the rational valuation $v$ is
$c\ord_P$, use a uniformizer of $P$ and separating residue
coordinates.  The extension has weights $(c,1)$ on that
uniformizer and $t$; the wedge of their logarithmic differentials
has log order zero, while the remaining coefficient has value
$cA_S(\ord_P)=A_S(v)$.  This proves the coefficient in
\eqref{eq:F-rees-log-order}.  There are no other divisor
coefficients, since off $E$ the form is the product of the
divisor-free $\omega$ and $dt/t$.

Since $K_S$ is Cartier and $v$ has integral values, $A_S(v)$ is an
integer.  Therefore
\[
 \eta=t^{-A_S(v)}\omega\wedge\frac{dt}{t}
 \qquad\text{satisfies}\qquad \divisor_{\mathcal S}(\eta)=-E.
\]
This establishes that the ordinary canonical divisor of the
total space is Cartier.  At every codimension-one point of $E$,
the fibre is regular and the flat morphism to the line is smooth.
The residue of $\eta$ is consequently a rational top form on $E$
with neither zeros nor poles at any codimension-one point.
It trivializes the reflexive canonical sheaf of $E$.
In particular this establishes ordinary canonical Cartierness,
not just Cartierness after adding a boundary.

Give $t$ weight $-1$ and $x_i t^{-a_i}$ weight $a_i$.  The residue
has weight $A_S(v)$ and the same $G$-character as $\omega$.
If an old torus is present, choose homogeneous generators;
$t$ is invariant for that torus, so its canonical character is
unchanged.  In the cone case such homogeneous generators have
common zero only at the vertex: their common zero set is
invariant and is isolated there, whereas the closure of every
positive-grading orbit contains the vertex.

Units of a positively graded affine domain with degree-zero part
$\C$ are constant.  Hence any two divisor-free canonical
generators on $E$ differ by a scalar.  The full fine grading
torus therefore acts on the generator by a character $\kappa$,
independent of the rational regrading used to construct it.
For all nearby rational $L$, its evaluation is
$L(A_S(w))$ by \eqref{eq:F-discrepancy-approximation} and the
residue computation.  Equality on this open set of rational
functionals identifies the two lattice weights and proves
\eqref{eq:F-canonical-weight}, including its equivariant version.

Finally suppose a nonidentity element of $G$ acquired a fixed
locus of codimension at most one on $E$.  At a general point of
that locus the fibre, and hence the morphism $\mathcal S\to\mathbb A^1$,
is smooth.  The fixed locus of the cyclic subgroup generated by
that element is smooth there.  Its tangent space is the invariant
tangent subspace, and it maps surjectively to the tangent space
of the line: averaging a lift of the base tangent vector gives
an invariant lift.  Thus this fixed locus is smooth over the
line and has the same fibre codimension on nearby nonzero
fibres.  This contradicts quasi-\'{e}taleness of the original
effective action.  The action on $S^0$ is therefore
quasi-\'{e}tale.
\end{proof}

\begin{corollary}[Lc equality and principal boundaries]
\label{cor:F-lc-degeneration}
If $w$ is $h$-lc, then $(S^0,\divisor(\bar h))$ is lc and
$\omega^0$ has the same fine torus weight and $G$-character as
$\bar h$.  Consequently $\omega^0/\bar h$ is invariant, and the
quotient cone has the log-form data of Section~\ref{sec:D}.
If $\Delta$ is any effective rational sum of principal divisors
near $x$, with $(S,\Delta)$ lc and
$A_S(w)=w(\Delta)$, its initial boundary $\Delta^0$ is lc on
$S^0$.  All these assertions also hold with specified additional
equivariant gradings.
\end{corollary}

\begin{proof}
Include the finitely many boundary equations in the presentation.
Their orders and discrepancy equality persist exactly under
rational regrading.  In \eqref{eq:F-rees}, each equation $f$ has
homogenization $f t^{-v(f)}$, which is a nonzerodivisor and has
no component $E$; it restricts to $\bar f$ on $E$.  In particular
the initial boundary is rationally Cartier.  The morphism
$\mathcal S\to S\times\mathbb A^1$ is birational.  Comparing the
coefficient at its only vertical prime by
\eqref{eq:F-rees-log-order} shows that
\[
 (\mathcal S,E+\Delta_{\mathcal S})
 \longrightarrow (S\times\mathbb A^1,
                       \Delta\times\mathbb A^1+S\times\{0\})
\]
is crepant precisely when $A_S(v)=v(\Delta)$.
Its source is therefore lc, and lc adjunction gives the stated
central pair.  At codimension-one points of the normal Cartier
fibre the morphism is smooth, so the different is exactly the
ordinary restriction of the homogenized boundary.
Apply this to $\Delta=H$.  Formula~\eqref{eq:F-canonical-weight}
and injectivity of evaluation on $\Lambda_w$ identify the weight
of $\omega^0$ with that of $\bar h$.  Quasi-\'{e}taleness and the
finite-map discrepancy formula descend klt, lc, and the invariant
log form to the quotient.  A suitable power of $\omega^0$
descends to a divisor-free pluricanonical form.
\end{proof}

\subsection{An exact bridge for inequalities}

\begin{proposition}[Sign transfer]\label{prop:F-bridge}
Suppose $w$ is common $h$-lc and $w(f_j)=0$ for every $j$.
Let $g=U/h^b$, where $0\ne U\in\mathcal O_{S,x}$,
$b\in\N_{>0}$, and $w(g)=0$.  Let $\mathcal Q$ be the actual
common slice and $\mathcal Q^0$ the common slice defined by
$\bar h$ and the $\bar f_j$ on $S^0$.  Then
\begin{equation}\label{eq:F-bridge}
 q(g)\le0\quad(q\in\mathcal Q)
 \quad\Longleftrightarrow\quad
 q^0(\bar g)\le0\quad(q^0\in\mathcal Q^0).
\end{equation}
One may add any finite list of tests of value zero at $w$.
There is no bound on the degree or the cost of $g$ in this statement.
\end{proposition}

\begin{proof}
Use Lemma~\ref{lem:E-slice-boundary} to encode the common slice.
Choose positive rational $\theta_j$ with $\sum_j\theta_j<1$ and
sufficiently general scalars $\lambda_j$, and set
\begin{equation}\label{eq:F-test-boundary}
 P_j=h^{m_j}+\lambda_jF_j,
 \qquad
 \Psi=(1-\textstyle\sum_j\theta_j)H+
                \sum_j\frac{\theta_j}{m_j}\divisor(P_j).
\end{equation}
The lemma says that $\Psi$ is effective and lc, with exactly the
common rays as its nonzero zero-discrepancy rays. It also gives
$q(P_j)=m_jq(h)$ on every common ray. These two properties let us
express a sign inequality as log canonicity of a small boundary
perturbation.

Assume the left side of \eqref{eq:F-bridge}.  For sufficiently
small positive rational $\delta$ the divisor
\begin{equation}\label{eq:F-perturbed-boundary}
 \Psi_\delta=\Psi+\delta(\divisor(U)-bH)
\end{equation}
is effective and lc near $x$.  Here is the exact reason for
smallness sufficing.  The coefficient of $H$ in
\eqref{eq:F-test-boundary} has a positive reserve, which preserves
effectivity.  On a fixed log resolution including $U$, every
coefficient below one stays below one for small $\delta$; for
a coefficient-one prime the discrepancy change is
$-\delta q(g)\ge0$.  Checking these finitely many components
checks lc on the entire resolution.  Restricting the germ if
necessary discards centres whose closed image does not contain
$x$.  No uniform choice of $\delta$ in $U$ is needed.

All equations in \eqref{eq:F-perturbed-boundary} have their
expected orders at $w$, and
$A_S(w)=w(\Psi_\delta)$ because $w(g)=w(f_j)=0$.
Corollary~\ref{cor:F-lc-degeneration} gives an lc initial boundary.
For every $q^0\in\mathcal Q^0$,
\[
 q^0(\overline{P_j})\ge m_jq^0(\bar h).
\]
Consequently
\[
 0\le A_{S^0,\Psi_\delta^0}(q^0)
       \le-\delta q^0(\bar g),
\]
which proves the right side of \eqref{eq:F-bridge}.

Conversely assume that right side.  Carry out exactly the same
boundary construction on $S^0$, choosing the $\lambda_j$ general
for both the actual and central constructions.  The intersection
of these finitely many nonempty open conditions is nonempty.
The assumed sign and a log resolution of the central equations
give a small rational $\delta>0$ for which $\Psi_\delta^0$ is
effective and lc near $o$.  Choose an integral rational regrading
that preserves all the equations involved.  Homogenize them in
the Rees family.  Their rational linear combination
$\Delta_{\mathcal S}$ is an effective rationally Cartier
horizontal boundary, disjoint in components from $E$, and its
restriction is $\Psi_\delta^0$.

The total space is normal, $E$ is normal and Cartier, and its
ordinary canonical divisor is Cartier by
Proposition~\ref{prop:F-canonical-degeneration}.  Thus
$K_{\mathcal S}+E+\Delta_{\mathcal S}$ is $\Q$-Cartier.
The different is the ordinary restriction, as checked above.
Inversion of adjunction for log canonicity
\citep[Theorem]{Kawakita} now makes
$(\mathcal S,E+\Delta_{\mathcal S})$ lc near $o$.
The section with all rescaled coordinates $x_i t^{-a_i}$ zero
specializes to $o$, so this neighbourhood meets it over a
nonzero parameter.  Off $E$ the family and its horizontal
boundary are the product with $(S,\Psi_\delta)$.  Hence the latter
pair is lc near $x$.
For a common actual ray, the initial boundary encoding has zero
discrepancy, and the last perturbation subtracts $\delta q(g)$.
It follows that $q(g)\le0$, proving the converse.

The resolution computations give inequalities on every divisorial
ray, and continuity on the monomial cones of the lc skeleton
gives them on the full quasi-monomial slices.  Additional finite
tests are treated in the same boundary and presentation.
\end{proof}

\begin{corollary}[Exact tests and pure weights]
\label{cor:F-exact-tests}
Suppose $w$ is common $h$-lc. If every common actual valuation
satisfies $q(f_j)=0$, then every common cone valuation satisfies
$q^0(\bar f_j)=0$.
The symbols $\bar f_j$ are invariant under $G$ and $\mathbb T_w$.
Every nonzero pure weight valuation of the cone is $h$-lc and,
after normalization by $\bar h=1$, is common.  Pure weights on
the boundary of the dual weight cone are allowed under the
closed-centre convention.
\end{corollary}

\begin{proof}
Apply Proposition~\ref{prop:F-bridge} to $g=f_j$, and combine its
upper bound with the defining lower bound.  The equality
$w(f_j)=0$ and injectivity of the fine grading imply that
$\bar f_j$ has weight zero.  The invariant log form
$\omega^0/\bar h$ has weight zero, so the pure-weight discrepancy
computation of Section~\ref{sec:D} gives lc equality for every
pure weight.  The test ratios then all have value zero.  Klt
implies that $\bar h$ is strictly positive on every nonzero
such weight, so the stated normalization is possible.
\end{proof}

\subsection{Compression of two gradings}

A second optimization takes place on the first cone.  To iterate the
construction on the original germ, we must realize these two successive
gradings by one actual valuation there.  The canonical-character
calculation will prove its lc equality, while the exact bridge keeps the
previous order constraints.

\begin{lemma}[Compression]\label{lem:F-compression}
Suppose $w$ is as above, is $h$-lc, and satisfies $w(f_j)=0$ for
every $j$. Let $v$ be a centred
$G\times\mathbb T_w$-invariant, $\bar h$-lc quasi-monomial
valuation on $S^0$, with normal finitely generated klt graded
algebra $R^{00}$.  Suppose also $v(\bar f_j)=0$.
The algebra $R^{00}$ has its joint grading by the old torus and
the $v$-grading.  For arbitrarily small positive $\delta$ outside
a countable exceptional set, it is the associated graded of an
actual centred quasi-monomial valuation $w_\delta$ on $S$, whose
grading evaluates a joint degree $(a,b)$ as $a+\delta b$.
The valuation $w_\delta$ is $h$-lc, is exact on all $f_j$, and,
after $h=1$ normalization, tends multiplicatively to $w$
uniformly on all nonzero regular functions.  If $v$ is nonzero
on an old degree-zero rational function, the rational rank of
$w_\delta$ is strictly larger than that of $w$.
\end{lemma}

\begin{proof}
Choose old-homogeneous elements of $\breve R$ whose $v$-symbols
generate $R^{00}$ and lift them to regular functions $x_i$ on
$S$.  Enlarge the list to generate the first graded algebra and
to include $h,F_j$ and every other specified equation.
Use $G$-eigenvectors throughout.  Put
$\alpha_i=w(x_i)$ and $\beta_i=v(\bar x_i)>0$.
Both successive subductions of
Lemma~\ref{lem:F-presentation} give
\begin{equation}\label{eq:F-iterated-initial}
 R^{00}=\C[X]/\operatorname{in}_\beta
                                  (\operatorname{in}_\alpha I).
\end{equation}
Completion of the first ideal does not change its initial ideal
for positive $\beta$.

Choose finitely many generators of the prime ideal on the right
that are iterated initials of elements of $I$.  This is possible
by first choosing $\alpha$-homogeneous elements of
$\operatorname{in}_\alpha I$ with the desired $\beta$-initials;
such elements lift with their prescribed $\alpha$-initial.
For each of these finitely many series, positivity and finiteness
of the relevant leading faces imply that its iterated initial
is its $(\alpha+\delta\beta)$-initial for all sufficiently small
$\delta>0$.  Hence
\[
 \operatorname{in}_\beta\operatorname{in}_\alpha I
           \subset\operatorname{in}_{\alpha+\delta\beta}I.
\]
The first ideal is prime and both quotients have dimension $n$.
They are equal.  Lemma~\ref{lem:F-presentation} supplies the
actual multiplicative valuation $w_\delta$ and its
quasi-monomiality.  Avoiding countably many values of $\delta$
makes $(a,b)\mapsto a+\delta b$ injective on the joint lattice.

We now prove the lc equality, after constructing the valuation.
The first degeneration has canonical character equal to the
weight of $\bar h$.  Apply
Proposition~\ref{prop:F-canonical-degeneration} to $v$ on that
cone, keeping the old torus action.  Since $v$ is $\bar h$-lc,
the resulting canonical generator has exactly the joint
character of $\bar{\bar h}$.  Independently apply the same
proposition to the actual valuation $w_\delta$, whose normal
finitely generated graded algebra has just been identified with
$R^{00}$.  Its discrepancy is the evaluation of this canonical
character under compression.  Therefore
\begin{equation}\label{eq:F-compressed-lc}
 A_S(w_\delta)
   =\langle w_\delta,\operatorname{wt}(\bar{\bar h})\rangle
   =w_\delta(h).
\end{equation}
This uses no prior lc assumption about $w_\delta$.
The included coordinates identify the specified symbols, and
the joint degree of each $\bar{\bar f_j}$ is zero; hence
$w_\delta(f_j)=0$.

For the uniform comparison, if
$(1-\varepsilon)\alpha_i\le\alpha_i+\delta\beta_i
\le(1+\varepsilon)\alpha_i$ for every $i$, the same inequalities
hold for every monomial weight.  The threshold-ideal description,
including its intersection with the actual local ring, gives
\[
 (1-\varepsilon)w(f)\le w_\delta(f)
                  \le(1+\varepsilon)w(f)
 \qquad(0\ne f\in\mathcal O_{S,x}).
\]
Normalization by $h$ changes the factors by another quantity
tending to one.  Finally the joint lattice surjects onto the old
lattice.  A degree-zero rational function with nonzero $v$-value
provides a nonzero vector in its kernel.  Injective compression
therefore strictly increases rational rank.
\end{proof}

\subsection{Balanced cones satisfying the angular order condition}

\begin{proposition}[Rank refinement]\label{prop:F-refinement}
Consider a sequence of the preceding data in dimensions at most
$n_0$.  Suppose that the common actual slices satisfy
$q(f_j)=0$ identically, that $w$ is a centred point of the slice
with the stipulated normal finitely generated klt graded algebra,
and that the cones satisfy the following balance condition with
one constant $C$ along the sequence:
\begin{equation}\label{eq:F-balance}
 q^0(s)\le Cw^0(s)
 \quad\begin{gathered}
 0\ne s\in\breve R\text{ fine-homogeneous},\\[-2pt]
 q^0\text{ common},\quad q^0(\bar h)=w(h)=1.
 \end{gathered}
\end{equation}
After passage to a subsequence, the $w$ can be replaced by $w'$
with all these properties, with a new fixed balance constant,
and with
\begin{equation}\label{eq:F-refinement-comparison}
 w'(f)=(1+o(1))w(f)
 \qquad(0\ne f\in\mathcal O_{S,x}),
\end{equation}
uniformly in $f$ in the normalization $w'(h)=w(h)=1$.
On the angular model $V_{\mathbf u}$ of the new cone, the
notation of Section~\ref{sec:D} then satisfies
\begin{equation}\label{eq:F-angular-order}
 T(\Gamma)\le(1+\eta_i)D(T),\qquad \eta_i\longrightarrow0,
\end{equation}
for every effective $\Q$-divisor $\Gamma\sim_\Q D$ and every
primitive divisorial $T$ such that
$a(T)=0$ and $T(\bar f_j)\ge0$ for all $j$.
This is condition \textup{(C-$\alpha$)} of Section~\ref{sec:C}.
The initial symbols here refer to $w'$.

The balance hypothesis holds when $w$ is the exact minimizer
on the common slice for the height $q(h)$.
\end{proposition}

\begin{proof}
The last assertion is the central domination estimate in
Corollary~\ref{cor:E-domination}; equivariance follows from the
uniqueness of the slice minimizer.  At every stage,
Corollary~\ref{cor:F-exact-tests} makes all common central test
values zero and makes every pure weight common.  The $\bar f_j$
are genuine functions of the angular field $\mathcal K$.
At a height-minimizing Gauss extension of $T$, regularity of
$\bar F_j$ gives
\begin{equation}\label{eq:F-angular-test-cost}
 T(\bar f_j)\ge-m_jD(T).
\end{equation}

Suppose \eqref{eq:F-angular-order} fails along a subsequence.
There are a fixed rational $c>0$, primitive witnesses $T$, and
effective divisors
\[
 \Gamma=D+\frac1m\divisor(\phi),\qquad m\in\N_{>0},
 \quad\phi\in\mathcal K^*,
\]
such that $a(T)=0$, $T(\bar f_j)\ge0$, and
\begin{equation}\label{eq:F-refinement-witness}
 T(\Gamma)-D(T)\ge2cD(T).
\end{equation}
Here $D(T)>0$ because $a+\mathbf D$ is klt.
After clearing denominators, the ray-ring description in
Section~\ref{sec:D} makes
$G_0=\bar h^m\phi$ a regular homogeneous function upstairs.
Choose a rational Gauss extension minimizing $\bar h$ at $T$.
It is $\bar h$-lc, and its lift to the finite quasi-\'{e}tale
cover cone remains lc.  Its centre has the vertex in its closed
image: this holds downstairs for a Gauss valuation, and finiteness
and the unique point over the vertex give it upstairs.
It is common, and \eqref{eq:F-refinement-witness} reads
\[
 \frac{g(q)}{q(\bar h)}\ge2c,
 \qquad g(q)=\frac1m q(G_0)-q(\bar h).
\]

Apply the one-sided tilt of Lemma~\ref{lem:E-one-sided-tilt} to
this cone.  Its exact minimizer $v$, normalized by $v(\bar h)=1$,
is centred, $G\times\mathbb T_w$-invariant, and has normal
finitely generated klt associated graded.  Localization at the
vertex does not change this graded algebra: a unit has constant
initial term.  The valuation $v$ is common, hence exact on the
old tests, and the tilt gives
\[
 v(\phi)>0,\qquad b(v)\ge1.
\]
The tilted height $b$ is bounded above by a fixed constant on
both the old and the new common slice, with constants depending
only on $n_0$ and the fixed $c$ at this stage.
Central domination therefore gives, on the new cone,
\begin{equation}\label{eq:F-tilt-balance}
 q^{00}(s)\le C_1v^0(s)
\end{equation}
for common $q^{00}$ normalized by $\bar{\bar h}=1$.
The joint canonical character equals that of $\bar{\bar h}$ by
Proposition~\ref{prop:F-canonical-degeneration}, including its
old-torus assertion.

Compress these two gradings by Lemma~\ref{lem:F-compression}.
Since $\phi$ has old weight zero and $v(\phi)>0$, rational rank
strictly increases.  The resulting actual valuation is $h$-lc,
exact on the tests, and can be chosen arbitrarily close to $w$
uniformly on every regular function.  We check balance, which
is necessary to repeat the argument.  A joint-homogeneous $s$
has an old-homogeneous lift $\widetilde s\in\breve R$; the
$v$-filtration splits into old weight spaces.  If its joint
degree is $(a,b)$, then
\[
 b=v(\widetilde s)\le Ca
\]
by \eqref{eq:F-balance}, since $v$ is common.  In the
normalization $w(h)=v(\bar h)=1$, the compressed degree is
$(a+\delta b)/(1+\delta)$.  Thus
\[
 v^0(s)=b\le C(1+\delta)w_\delta^0(s).
\]
Together with \eqref{eq:F-tilt-balance} this is the next
balance estimate, with a fixed constant.  A generic compression
makes every fine-homogeneous element joint-homogeneous, so this
checks precisely all elements required in \eqref{eq:F-balance}.

If the angular order condition still fails, pass to a subsequence
with a fixed positive defect and repeat.  Each repetition raises
rational rank, which is at most the dimension.  There can be at
most $n_0$ repetitions.  At a final stage the supremum of the
positive normalized defects must tend to zero.  This supremum is
finite: a height-minimizing Gauss witness and
\eqref{eq:F-balance}, applied to $G_0$ of grading value $m$,
bound $T(\Gamma)/D(T)$ by the balance constant.
Taking these suprema, or slightly larger positive numbers, gives
the single sequence $\eta_i$ in \eqref{eq:F-angular-order},
uniformly for all its displayed tests and divisors.
At each of the boundedly many compressions choose the
all-functions multiplicative error to tend to zero; their product
still tends to one and proves \eqref{eq:F-refinement-comparison}.
All exact test equalities, equivariance, canonical characters,
and quasi-\'{e}taleness have been preserved at every step.
\end{proof}
